% from ega4-18-fr.tex
\oldpage[IV-4]{189}

et de $\mathfrak c$ dans cet anneau quotient, montre que $\mathfrak c\subset g_1B$. Soit $\mathfrak d$ l'idéal de $B$ formé des $x\in B$ tels que $g_1x\in\mathfrak c$; on a donc $\mathfrak c=g_1\mathfrak d$; mais comme par hypothèse $g_1$ est régulier dans $B/\mathfrak c$, on a nécessairement $\mathfrak d=\mathfrak c$, donc $\mathfrak c=g_1\mathfrak c$. Comme $\mathfrak c$ est de type fini et que $g_1$ est contenu dans le radical de $B$, le lemme de Nakayama montre que $\mathfrak c=0$.

Démontrons enfin l'assertion (iii) de (19.1.6). En vertu de (i) et de (19.1.2), il suffit de voir que si $\mathfrak R$ est régulier, il en est de même de $\mathfrak R'$. Notons que $\mathfrak J$ est régulier (19.1.2), et raisonnons par récurrence sur le rang $q$ du $(A/\mathfrak J)$-module libre $\mathfrak J/\mathfrak J^2$. Si $q=0$, on a $\mathfrak J=0$ par le lemme de Nakayama et l'assertion est triviale. Comme $A$ et $A/\mathfrak J$ sont réguliers, on sait ($0$, 17.1.9) que $\mathfrak J$ est engendré par une partie à $q$ éléments d'un système régulier de paramètres de $A$, donc, si $f$ est un des éléments de ce système de générateurs de $\mathfrak J$, $A_1=A/fA$ est régulier et $f\notin\mathfrak m^2$ ($0$, 17.1.8). Soient $\mathfrak J_1$, $\mathfrak R_1$ les images canoniques de $\mathfrak J$, $\mathfrak R$ dans $A_1$; on a $\mathfrak J_1\subset\mathfrak R_1$, $A_1/\mathfrak J_1=A/\mathfrak J$, donc $A_1/\mathfrak J_1$ est régulier et par suite $\mathfrak J_1$ est un idéal régulier (19.1.2). Montrons que $\mathfrak R_1$ est un idéal régulier dans $A_1$; comme $\mathfrak R$ est un idéal régulier dans $A$, il suffit de montrer que $f$ fait partie d'un système régulier de générateurs de $\mathfrak R$, et pour cela (16.9.5) il suffit de montrer que l'image de $f$ dans $\mathfrak R/(\mathfrak R^2+\mathfrak m\mathfrak R)=\mathfrak R/\mathfrak m\mathfrak R$ fait partie d'une base de cet $(A/\mathfrak m)$-espace vectoriel, autrement dit que $f\notin\mathfrak m\mathfrak R$, ce qui résulte en effet de $f\notin\mathfrak m^2$. Alors, comme $\mathfrak J_1/\mathfrak J_1^2$ est un $(A_1/\mathfrak J_1)$-module libre de rang $q-1$, l'hypothèse de récurrence montre que $\mathfrak R'_1=\mathfrak R_1/\mathfrak J_1$ est un idéal régulier de $A_1/\mathfrak J_1$, et comme $\mathfrak R'$ s'identifie à $\mathfrak R'_1$ dans l'isomorphisme canonique entre $A/\mathfrak J$ et $A_1/\mathfrak J_1$, $\mathfrak R'$ est régulier. C.Q.F.D.

\begin{remarks}[19.1.7]
\label{IV.19.1.7}
\begin{enumerate}
  \item[(i)] Nous ignorons si le composé de deux immersions quasi-régulières est quasi-régulière.
  \item[(ii)] Dans un espace annelé en anneaux locaux $X$, pour tout Idéal $\mathscr J$ de $\sh{O}_X$, $\sh{O}_X/\mathscr J$ est de type fini et a donc un support fermé dans $X$ ($0_{\mathrm I}$, 5.2.2). On peut par suite définir comme dans (I, 4.1.3 et 4.2.1) les notions de sous-espace annelé (localement fermé) de $X$ défini par un Idéal $\mathscr J$ d'un faisceau $\sh{O}_U$, où $U$ est un ouvert de $X$, et la notion d'immersion. Les définitions d'immersions quasi-régulières et régulières, de codimension transversale s'appliquent alors sans modification, et les résultats de (19.1.4, (i) à (iv)) sont encore valables, en supposant dans la seconde assertion de (i) et dans (iv) que le faisceau $\sh{O}_X$ soit cohérent et les anneaux locaux $\sh{O}_{X,x}$ noethériens; quant à (ii), il faut \emph{définir} $Y'$ comme le sous-espace annelé défini par le $\sh{O}_{X'}$-Idéal $g^*(\mathscr J)$; les démonstrations sont inchangées.
\end{enumerate}
\end{remarks}

\begin{env}[19.1.8]
\label{IV.19.1.8}
Nous avons déjà signalé (notamment dans (16.9.6, (ii))) que dans les anneaux non noethériens, la notion de « suite régulière » d'éléments de l'anneau ne possède pas les propriétés qui la rendraient utilisable. De récentes recherches semblent montrer que la notion qui doit dans ce cas se substituer à celle de suite régulière est celle de suite $\mathbf f=(f_1,\ldots,f_n)$ ayant la propriété que le complexe de l'algèbre extérieure $\mathrm K_0(\mathbf f,M)$ (III, 1.1.3) ait \emph{son homologie nulle en degrés $>0$} (cf. (III, 1.1.4)); voir en particulier A. Grothendieck, Séminaire de géométrie algébrique, 1966, à paraître prochainement.
\end{env}

\subsection{Immersions transversalement régulières.}
\EGAReaderTarget{ega4.c04.subsec.19-2}{subsection:IV.19.2}

